\scanpage{139}
\mspage{359}

to a unique $u'$ such that $u'(\sigma_\infty)=\sigma_\infty$ i.e. a $u'$ which commutes with $\sigma_\infty$. But in \unsure{the symmetric notation|the symmetric notations} we shall take
\[
u''=\operatorname{int}(\alpha^{-1})u
\]
corresponding to
\[
\left\{\begin{array}{l}
g''=\alpha^{-1}g\rho(\alpha)=1\\
h''=\alpha^{-1}h\sigma_\infty(\alpha)=(\alpha^{-1}\beta)\sigma_\infty(\alpha^{-1}\beta)^{-1}\\
f''=f\alpha
\end{array}\right.
\qquad (17)
\]
but $u''$ will be \inserted{only $\widehat{\widetilde{\Pi}}_{0,3}$-}conjugate \deleted{\unsure{by}} of $\unsure{u''|u}$ -- it will be \inserted{$\widehat{\Pi}_{0,3}$-}conjugate to an \deleted{element of} \inserted{\unsure{automorphism}} which commutes with $\rho$ iff $g$ is $\Pi_{0,3}$-$\rho$-conjugate to $1$ i.e. iff $\alpha\in\widehat{\Pi}_{0,3}$ i.e. iff $\mu\equiv1\ (3)$ or again $\nu\equiv0\ (3)$.

\deleted{As for the} The calculations of the preceding \S\S\ concerned the case where $h=1$, i.e. $u$ commutes with $\sigma_\infty$ -- so \illegible{1} there remain the quantities $g$ and $f$, \deleted{we see moreover} \inserted{plus the multiplier $\mu$}
\deleted{\inserted{the multiplier} for $\mu$, \unsure{or for} $\nu=\frac{\mu-1}{2}$, is determined unambiguously by}
-- once $\mu$ is known i.e. $\nu$ known, $g$ and $f$ determine one another by (11) ($f$ remaining an element of $L_0\backslash\widehat{\Pi}_{0,3}$).

\margin{We shall see that $\mu$ is already determined by $f$, \unsure{at least modulo} sign, cf. below.}

For fixed $f$, the equation in $g$ can be interpreted as an equation in $f$, or better, introducing
\[
\varphi=f\sigma_\infty(f)^{-1}
\qquad (12)
\]

\margin{This notation is \unsure{useful|valid} only if $\beta=1$ \illegible{2}}
\margin{it is essentially the element denoted $h_0$ previously}
\margin{138}

\scanpage{140}

an ``equation in $\varphi$'' --
\[
(\ell_0^\nu\varphi)\rho(\ell_0^\nu\varphi)\rho^2(\ell_0^\nu\varphi)=1,
\qquad (13)
\]
or again, putting
\[
\boxed{\xi=\ell_0^\nu\varphi\quad\text{i.e. }\varphi=\ell_0^{-\nu}\xi}
\qquad (14)
\]
we can interpret it as ``equation in $\xi$''
\[
\xi\rho(\xi)\rho^2(\xi)=1.
\qquad (15)
\]
Thus, the loop automorphisms of $\widehat{\Pi}_{0,3}$ which commute with $\sigma_\infty$, \inserted{with multiplier $\mu=2\nu+1$,} correspond to the pairs of solutions $(\xi,\varphi)$ of the equations
\[
\boxed{\varphi\sigma_\infty(\varphi)=1\qquad
\xi\rho(\xi)\rho^2(\xi)=1}
\qquad (16)
\]

\margin{We \unsure{shall find|find} that \unsure{both} $\xi$ and $\varphi$ already \unsure{determine} $\mu$, hence $\nu$.}

satisfying in addition the compatibility condition (14) (particularly simple \deleted{\unsure{(17)}} if $\mu=1$ i.e. \deleted{$\mu=$} $\nu=0$, since then $\varphi=\xi$ !). We recover $g$ (hence $u$) \deleted{by the} as a function
\deleted{Conversely, if we \illegible{3}}
of $(\varphi,\xi)$ by the loop formula (11), which here takes the form
\[
g=\sigma_\infty(f)^{-1}\ell_1^\nu\rho(f)
=(\ell_0^\nu f)^{-1}\bigl(\underbrace{\ell_0^\nu\varphi}_{\xi}\bigr)\rho(\ell_0^\nu f)
\qquad (17)
\]
$f\in\widehat{\Pi}_{0,3}$ being uniquely determined in terms of $\varphi$ by formula (12).

Symmetrically, we are going to look at the loop automorphisms $u$ of $\Pi_{0,3}$ which com-

\margin{139}

\scanpage{141}
\mspage{360}

mute with $\rho$, i.e. those for which we have $g=1$, hence described in terms of $h,f$,
\[
u(\sigma_\infty)=h\sigma_\infty,\quad u(\rho)=1_{\deleted{\unsure{\rho}}}\qquad (18)
\]
\[
u(\varepsilon_0)=\operatorname{int}(f^{-1})(\varepsilon_0^\mu)
\qquad (19)
\]
this last relation taking the form (11)
\[
\sigma_\infty(h)=\sigma_\infty(f)^{-1}\ell_1^\nu\rho(f),
\qquad (20)
\]
which can also be written, applying $\sigma_\infty$
\[
h=f^{-1}.\ell_0^\nu.(\sigma_\infty\rho)(f)\unsure{\underset{\cdot}{=}|=}
\qquad (21)
\]
``The equation in $h$'' (2) \unsure{the}
\[
h\sigma_\infty(h)=1
\qquad \unsure{(*)}
\]
is thus written as an ``equation in $f$'', \deleted{which} \inserted{with}
\[
h\sigma_\infty(h)=f^{-1}\ell_0^\nu
\underbrace{(\sigma_\infty\rho)(f)\sigma_\infty(f)^{-1}}_{\sigma_\infty(\underbrace{\rho(f)f^{-1}}_{\psi^{-1}})}
\ell_1^\nu\rho(f)\deleted{\illegible{1}}
\]
so \unsure{this} is written, in terms of $\psi=f\rho(f)^{-1}$
\[
\underbrace{\ell_0^\nu\sigma_\infty(\psi^{-1})}\ell_1^\nu
\underbrace{\rho(f)f^{-1}}_{\psi^{-1}}=1
\]
or (passing to the inverse)
\[
\psi\ell_1^{-\nu}\underbrace{\sigma_\infty(\psi)\ell_0^{-\nu}}_{\sigma_\infty(\psi\ell_1^{-\nu})}=1
\quad\text{or}\quad
\underbrace{\ell_0^{-\nu}\psi}_{\eta}
\underbrace{\ell_1^{-\nu}\sigma_\infty(\psi)}_{\sigma_\infty(\underbrace{\ell_0^{-\nu}\psi}_{\eta})}=1.
\]
\deleted{or}
Thus, putting
\[
\psi\overset{\mathrm{dfn}}{=}f\rho(f)^{-1}
\qquad (22)
\]
\[
\boxed{\eta=\ell_0^{-\nu}\psi\quad\text{i.e. }\psi=\ell_0^\nu\eta}
\qquad (23)
\]
, we find that the loop automorphisms $u$ of $\widehat{\Pi}_{0,3}$ which commute with $\rho$ correspond to the triples $(\eta,\psi,\nu)$, $\eta,\psi\in\widehat{\Pi}_{0,3}$, $\nu\in\widehat{\mathbb Z}$ (with $2\nu+1\in\widehat{\mathbb Z}^{*}$), satisfying

\margin{140}

\scanpage{142}

\[
\boxed{\psi\rho(\psi)\rho^2(\psi)=1,\qquad\eta\sigma_\infty(\eta)=1},
\qquad (24)
\]
\margin{We \unsure{shall find} \illegible{1} \unsure{again}, \unsure{by} \unsure{both} $\psi$ and $\eta$}

and in addition the compatibility relation (23) -- particularly simple if $\mu=1$ i.e. if $\nu=0$, for then it means $\psi=\eta$ ! But in addition to these relations, we must assume that the solution $\psi$ of equation (24) ``in $\psi$'' \unsure{belongs|belong} to the principal family, i.e. that, when we write $\psi$ in parametric form
\[
\psi=f\rho(f)^{-1}
\qquad (22)
\]
with $f\in\widehat{\widetilde{\Pi}}_{0,3}=\widehat{\Pi}_{0,3}.\{1,\tau_\infty\}$, we indeed have $f\in\widehat{\Pi}_{0,3}$ (hence that the image of $f$ in $\widehat{\widetilde{\Pi}}_{0,3}/\widehat{\Pi}_{0,3}\simeq\{1,\tau_\infty\}\simeq\mathbb Z/2$ is zero\dots).

Having thus determined $f$ as a \uline{function} of $\psi$, \unsure{determine|determines} $h$ by (21), whence $u$ \deleted{\illegible{3}}.

\uline{Remarks} 1) If this additional condition on $\psi$ is not satisfied, the \deleted{\unsure{determination}} formula (21) still makes sense and defines $h\in\widehat{\Pi}_{0,3}$, satisfying ``the equation in $h$'' $h\sigma_\infty(h)=1$ hence defining an automorphism $u$ of $\widehat{\Pi}_{0,3}$, commuting with $\rho$.

We shall then have
\[
u(\ell_0)=\operatorname{int}(f^{-1})(\ell_0^\mu)
\]
or again, putting
\[
f=\inserted{\tau_\infty}f_0\deleted{\tau_\infty}\qquad\text{with }f_0\in\widehat{\Pi}_{0,3}
\]
\[
u(\ell_0)=\operatorname{int}(f_0^{-1})(\tau_\infty(\ell_0^\mu))=\operatorname{int}(f_0^{-1})(\ell_0^{-\mu})
\]

\margin{141}

\scanpage{143}
\mspage{361}

i.e. $u$ is a loop automorphism, with multiplier $-\mu=-(2\nu+1)$.

2) Let us note that the indeterminacy of $f$ (for fixed $u$, i.e. fixed $g,h$) lies in
\[
 f\longmapsto f'=\ell_0^n f\;\deleted{\unsure{f}}
 \qquad(n\in\widehat{\mathbb Z})
 \tag{25}
\]
which gives
\[
 \left\{\begin{aligned}
 \varphi&\longmapsto\varphi'=\ell_0^n\varphi\ell_1^{-n},&
 \psi&\longmapsto\psi'=\ell_0^n\psi\ell_1^{-n},\\
 \xi&\longmapsto\xi'=\ell_0^n\xi\ell_1^{-n},&
 \eta&\longmapsto\eta'=\ell_0^n\eta\ell_1^{-n},
 \end{aligned}\right.
 \tag{26}
\]
\margin{\uline{N.B.} We similarly have
$\alpha_0\longmapsto\alpha_0'=\ell_0^n\alpha_0$
$\beta_0\longmapsto\beta_0'=\ell_0^n\beta_0$.}

(N.B. we have $\ell_1=\sigma_\infty(\ell_0)=\rho(\ell_0)$, so $\xi'$ and $\psi'$ are the $\rho$-transforms of $\xi,\psi$ by $\unsure{\ell_0^{-n}|\ell_0^n}$, and $\varphi',\eta'$ the $\sigma_\infty$-transforms of $\varphi,\eta$ by $\unsure{\ell_0^{-n}|\ell_0^n}$ --- in particular the equations (16) (in $\varphi$ and in $\xi$) and (24) (in $\psi$ and in $\eta$) are indeed invariant under these transformations, as are moreover the compatibility relations (14) resp. (23).

3) We have the following formulas \inserted{(rewriting of (11))}
\[
 \sigma_\infty(h)g
 = (\ell_0^\nu f)^{-1}\xi\rho(\ell_0^\nu f)
 = \sigma_\infty(f^{-1}\ell_0^\nu)\eta(f^{-1}\ell_0^\nu)^{-1}
 \tag{27}
\]
\deleted{(case $h=1$)} --- expressing $\inserted{\sigma_\infty(h)}g$ as the $\rho$-conjugate of $\xi$ by $(\ell_0^\nu f)^{-1}$, and symmetrically \deleted{in the case $g=1$,} as the $\sigma_\infty$-conjugate of $\eta$ by $\sigma_\infty(f^{-1}\ell_0^\nu)$. If $h=1$, it is the first expression that is interesting (implying the equation in $\xi$), for $g=1$, it is the second, which we can also rewrite \inserted{(applying $\sigma_\infty$ to the 2 sides)}
\[
 h=(f^{-1}\ell_0^\nu)(\eta^{-1})\sigma_\infty(f^{-1}\ell_0^\nu)^{-1}
 \tag{28}
\]
\margin{142}

\scanpage{144}

expressing $h$ as $\sigma_\infty$-transform of $\eta^{-1}$ by $f^{-1}\ell_0^\nu$. If we want to express it as \inserted{$\sigma_\infty$-}transform of $\eta$ (instead of $\inserted{\sigma_\infty(\eta)=}\eta^{-1}$), we write
\[
 \deleted{\unsure{\eta^{-1}=\eta^{-1}\eta\sigma_\infty(\eta^{-1})}}
\]
\[
 \eta^{-1}=\sigma_\infty(\eta)=\eta^{-1}.\eta.\sigma_\infty(\eta)
\]
whence
\[
 h=(f^{-1}\ell_0^\nu\eta^{-1})\eta\sigma_\infty(f^{-1}\ell_0^\nu\eta^{-1})^{-1}.
 \tag{28bis}
\]

These formulas allow us to relate, in the case $h=1$ the solutions of the \deleted{\unsure{two}} equations
\[
 g=\alpha\rho(\alpha)^{-1}
 \qquad\text{and}\qquad
 \xi=\alpha_0\rho(\alpha_0)^{-1}
 \tag{29}
\]
by
\[
 \alpha=(\ell_0^\nu f)^{-1}\alpha_0=f^{-1}\ell_0^{-\nu}\alpha_0,
 \tag{30}
\]
and in the case where $g=1$, the solutions of the two equations
\[
 h=\beta\sigma_\infty(\beta)^{-1},\qquad
 \eta=\beta_0\sigma_\infty(\beta_0)^{-1}
 \tag{31}
\]
by
\[
 \beta=(f^{-1}\ell_0^\nu\eta^{-1})\beta_0.
 \tag{32}
\]

The asymmetry between these two formulas ((32) being a little more complicated than (30) by the extra factor $\eta^{-1}$ in the coefficient of $\beta_0$) comes from the fact that $f$ was defined via the loop \unsure{formula|formulas} (6), first form, where $\ell_0$ is defined as the square of $\sigma_\infty\rho=\varepsilon_0$, and not of $\rho\sigma_\infty=\varepsilon_1$
\margin{143}